\chapter{Complexity — Levels}
%%% generated by the blueprint pipeline from TCSlib.Complexity.PolyHierarchy.Levels
\section{Overview}

Declarations of the Lean module \texttt{TCSlib.Complexity.PolyHierarchy.Levels}.

\section{Declarations}

\begin{theorem}[Preimage closure at the existential hierarchy level]
\lean{Complexity.preimage_mem_SigmaP}
\label{Complexity.preimage_mem_SigmaP}
\leanok
\uses{Complexity.PolyHierarchy.preimage_mem_altClass, Complexity.PolyTimeComputable, Complexity.SigmaP}
\difficulty{1}
Let $i\in\mathbb{N}$, let $L$ belong to the $i$th existential level of the polynomial
hierarchy, and let $f$ be a polynomial-time computable function from binary strings to
binary strings. Then $\{x:f(x)\in L\}$ belongs to the $i$th existential level of the
polynomial hierarchy.
\end{theorem}

\begin{theorem}[Preimage closure at the universal hierarchy level]
\lean{Complexity.preimage_mem_PiP}
\label{Complexity.preimage_mem_PiP}
\leanok
\uses{Complexity.PiP, Complexity.PiP_eq_altClass, Complexity.PolyHierarchy.preimage_mem_altClass, Complexity.PolyTimeComputable}
\difficulty{3}
Let $i\in\mathbb{N}$, let $L$ belong to the $i$th universal level of the polynomial
hierarchy, and let $f$ be a polynomial-time computable function from binary strings to
binary strings. Then $\{x:f(x)\in L\}$ belongs to the $i$th universal level of the
polynomial hierarchy.
\end{theorem}

\begin{definition}[Certificate split encoding]
\lean{Complexity.certSplit}
\label{Complexity.certSplit}
\leanok
\uses{Turing.pairEncode, Turing.solveSplit}
Let $C,c$ be nonnegative integers and let $w$ be a bit string. If the split-selection
procedure returns an index $i$ for the length of $w$, encode the prefix of $w$ of length
$i$ together with the remaining suffix using the self-delimiting pairing; if it returns
no index, output the empty bit string.
\end{definition}

\begin{theorem}[A solution of the padding equation]
\lean{Complexity.solveSplit_self}
\label{Complexity.solveSplit_self}
\leanok
\uses{Turing.solveSplit}
\difficulty{4}
For all natural numbers $C,c,n$, the least natural number $i\le n+C(n+1)^c$ satisfying
$i+C(i+1)^c=n+C(n+1)^c$ is $n$.
\end{theorem}

\begin{theorem}[Certificate split on concatenated strings]
\lean{Complexity.certSplit_append}
\label{Complexity.certSplit_append}
\leanok
\uses{Complexity.certSplit, Complexity.solveSplit_self, Turing.pairEncode}
\difficulty{1}
Let $C,c$ be natural numbers and let $x,u$ be finite bit strings. If $|u|=C(|x|+1)^c$,
then the certificate split of the concatenation $xu$ is the self-delimiting pair
encoding of $x$ and $u$.
\end{theorem}

\begin{theorem}[The first existential level]
\lean{Complexity.SigmaP_one}
\label{Complexity.SigmaP_one}
\leanok
\uses{Complexity.NP, Complexity.PolyTimeComputable, Complexity.SigmaP, Complexity.altQuant_true_succ, Complexity.certSplit, Complexity.certSplit_append, Complexity.pairFstD, Complexity.pairSndD, Complexity.polyTimeComputable_pairConcat, Complexity.preimage_mem_P, Turing.FinTM.computesFunInTime_splitSolve, Turing.pairEncode}
\difficulty{5}
The class $\Sigma_1^P$ equals $\mathsf{NP}$, the class of languages whose membership has
polynomial-length certificates verifiable in polynomial time.
\end{theorem}

\begin{theorem}[First universal polynomial hierarchy level]
\lean{Complexity.PiP_one}
\label{Complexity.PiP_one}
\leanok
\uses{Complexity.NP, Complexity.PiP, Complexity.SigmaP, Complexity.SigmaP_one, Complexity.coNP}
\difficulty{3}
The class of languages whose complements belong to the first existential alternating
level with polynomial-time verification is exactly the class of languages whose
complements are in $mathrm{NP}$; equivalently, $mathrm{PiP}(1)=mathrm{coNP}$.
\end{theorem}

\begin{theorem}[Successor-level characterization of the polynomial hierarchy]
\lean{Complexity.mem_SigmaP_succ_iff}
\label{Complexity.mem_SigmaP_succ_iff}
\leanok
\uses{Complexity.PiP, Complexity.PiP_eq_altClass, Complexity.PolyHierarchy.UnaryPT, Complexity.PolyHierarchy.mem_altClass_of_normal, Complexity.PolyHierarchy.mem_altClass_of_uniform, Complexity.PolyHierarchy.unaryPT_poly, Complexity.PolyTimeComputable.pairEncode, Complexity.SigmaP, Complexity.altQuant, Complexity.altQuant_true_succ, Complexity.mem_PiP_iff, Complexity.pairFstD, Complexity.pairFstD_pairEncode, Complexity.polyTimeComputable_id, Complexity.polyTimeComputable_pairFstD, Turing.length_pairEncode, Turing.pairEncode}
\difficulty{5}
For every natural number $i$ and every language $L$ of binary strings, $L$ belongs to
the $(i+1)$st existential level of the polynomial hierarchy if and only if there are
natural numbers $C,c$ and a language $L'$ in the $i$th universal level such that, for
every binary string $x$, membership $x\in L$ is equivalent to the existence of a binary
string $u$ of length exactly $C(|x|+1)^c$ for which the self-delimiting pair encoding
$\langle x,u\rangle$ belongs to $L'$.
\end{theorem}

\begin{theorem}[Characterization of the successor hierarchy class]
\lean{Complexity.mem_PiP_succ_iff}
\label{Complexity.mem_PiP_succ_iff}
\leanok
\uses{Complexity.PiP, Complexity.SigmaP, Complexity.compl_mem_PiP_iff, Complexity.compl_mem_SigmaP_iff, Complexity.mem_SigmaP_succ_iff, Turing.pairEncode}
\difficulty{4}
For every natural number $i$ and every language $L$ of finite binary strings,
$L\in\Pi_{i+1}^P$ if and only if there are natural numbers $C,c$ and a language
$L'\in\Sigma_i^P$ such that, for every binary string $x$, membership $x\in L$ is
equivalent to the condition that every binary string $u$ of length $C(|x|+1)^c$
satisfies $\langle x,u\rangle\in L'$.
\end{theorem}

\begin{theorem}[Closure under existential polynomial witnesses]
\lean{Complexity.mem_SigmaP_of_exists}
\label{Complexity.mem_SigmaP_of_exists}
\leanok
\uses{Complexity.PolyHierarchy.UnaryPT, Complexity.PolyHierarchy.UnaryPT.add, Complexity.PolyHierarchy.mem_altClass_of_normal, Complexity.PolyHierarchy.padDecode, Complexity.PolyHierarchy.padDecode_marker, Complexity.PolyHierarchy.polyTimeComputable_padDecode, Complexity.PolyHierarchy.unaryPT_const, Complexity.PolyHierarchy.unaryPT_poly, Complexity.PolyTimeComputable, Complexity.PolyTimeComputable.comp, Complexity.PolyTimeComputable.pairEncode, Complexity.SigmaP, Complexity.altQuant_true_succ, Complexity.pairFstD, Complexity.pairSndD, Complexity.polyTimeComputable_id, Complexity.polyTimeComputable_pairFstD, Complexity.polyTimeComputable_pairSndD, Complexity.qStep, Turing.length_pairEncode, Turing.pairEncode}
\difficulty{6}
Let $i$ be a nonzero natural number, and let $L,L'$ be languages of finite binary
strings with $L'\in\Sigma_i^P$. For any natural numbers $C,c$, if for every binary
string $x$, membership $x\in L$ is equivalent to the existence of a binary string $u$ of
length $C(|x|+1)^c$ such that $\langle x,u\rangle\in L'$, then $L\in\Sigma_i^P$.
\end{theorem}

\begin{theorem}[Closure under universal polynomial witnesses]
\lean{Complexity.mem_PiP_of_forall}
\label{Complexity.mem_PiP_of_forall}
\leanok
\uses{Complexity.PiP, Complexity.mem_SigmaP_of_exists, Turing.pairEncode}
\difficulty{3}
Let $i$ be a nonzero natural number, and let $L,L'$ be languages of finite binary
strings with $L'\in\Pi_i^P$. For any natural numbers $C,c$, if for every binary string
$x$, membership $x\in L$ is equivalent to membership of $\langle x,u\rangle$ in $L'$ for
every binary string $u$ of length $C(|x|+1)^c$, then $L\in\Pi_i^P$.
\end{theorem}

\begin{theorem}[Pi Level Inclusion in the Next Sigma Level]
\lean{Complexity.PiP_subset_SigmaP_succ}
\label{Complexity.PiP_subset_SigmaP_succ}
\leanok
\uses{Complexity.PiP, Complexity.SigmaP, Complexity.mem_SigmaP_succ_iff, Complexity.pairFstD, Complexity.pairFstD_pairEncode, Complexity.polyTimeComputable_pairFstD, Complexity.preimage_mem_PiP, Turing.pairEncode}
\difficulty{3}
For every $i\in\mathbb{N}$, every language in $\Pi_i^P$ belongs to $\Sigma_{i+1}^P$.
\end{theorem}

\begin{theorem}[Existential levels and the next universal level]
\lean{Complexity.SigmaP_subset_PiP_succ}
\label{Complexity.SigmaP_subset_PiP_succ}
\leanok
\uses{Complexity.PiP, Complexity.PiP_subset_SigmaP_succ, Complexity.SigmaP, Complexity.compl_mem_PiP_iff}
\difficulty{1}
For every natural number $i$, every language in $\Sigma_i^P$ belongs to $\Pi_{i+1}^P$,
the class of languages whose complements belong to $\Sigma_{i+1}^P$.
\end{theorem}

\begin{theorem}[Successive existential levels are nested]
\lean{Complexity.SigmaP_subset_SigmaP_succ}
\label{Complexity.SigmaP_subset_SigmaP_succ}
\leanok
\uses{Complexity.PiP_subset_SigmaP_succ, Complexity.PiP_zero, Complexity.SigmaP, Complexity.SigmaP_zero, Complexity.mem_SigmaP_succ_iff}
\difficulty{4}
For every natural number $i$, every language in $\Sigma_i^P$ belongs to
$\Sigma_{i+1}^P$.
\end{theorem}

\begin{theorem}[Inclusion of Consecutive Pi Levels]
\lean{Complexity.PiP_subset_PiP_succ}
\label{Complexity.PiP_subset_PiP_succ}
\leanok
\uses{Complexity.PiP, Complexity.SigmaP_subset_SigmaP_succ}
\difficulty{1}
For every $i\in\mathbb{N}$, every language in $\Pi_i^P$ belongs to $\Pi_{i+1}^P$.
\end{theorem}

\begin{theorem}[Monotonicity of the existential hierarchy]
\lean{Complexity.SigmaP_mono}
\label{Complexity.SigmaP_mono}
\leanok
\uses{Complexity.SigmaP, Complexity.SigmaP_subset_SigmaP_succ}
\difficulty{4}
For all natural numbers $i$ and $j$ with $i \le j$, every language in $\Sigma_i^P$
belongs to $\Sigma_j^P$.
\end{theorem}

\begin{theorem}[Monotonicity of the Complement Alternation Classes]
\lean{Complexity.PiP_mono}
\label{Complexity.PiP_mono}
\leanok
\uses{Complexity.PiP, Complexity.SigmaP_mono}
\difficulty{1}
For natural numbers $i\leq j$, every language in $\Pi_i^P$ also belongs to $\Pi_j^P$.
\end{theorem}

\begin{theorem}[Adjacent levels contain existential and universal classes]
\lean{Complexity.SigmaP_union_PiP_subset}
\label{Complexity.SigmaP_union_PiP_subset}
\leanok
\uses{Complexity.PiP, Complexity.PiP_subset_PiP_succ, Complexity.PiP_subset_SigmaP_succ, Complexity.SigmaP, Complexity.SigmaP_subset_PiP_succ, Complexity.SigmaP_subset_SigmaP_succ}
\difficulty{4}
For every natural number $i$, the union $\Sigma_i^P\cup\Pi_i^P$ is contained in
$\Sigma_{i+1}^P\cap\Pi_{i+1}^P$.
\end{theorem}

\begin{theorem}[Polynomial time lies in every existential level]
\lean{Complexity.P_subset_SigmaP}
\label{Complexity.P_subset_SigmaP}
\leanok
\uses{Complexity.P, Complexity.SigmaP, Complexity.SigmaP_mono, Complexity.SigmaP_zero}
\difficulty{3}
For every natural number $i$, every binary language decidable in polynomial time belongs
to the $i$th existential class $\Sigma_i^{\mathsf{P}}$: $\mathsf{P} \subseteq
\Sigma_i^{\mathsf{P}}$.
\end{theorem}

\begin{theorem}[Polynomial time lies in every universal level]
\lean{Complexity.P_subset_PiP}
\label{Complexity.P_subset_PiP}
\leanok
\uses{Complexity.P, Complexity.PiP, Complexity.PiP_mono, Complexity.PiP_zero}
\difficulty{1}
For every natural number $i$, every binary language decidable in polynomial time belongs
to the $i$th universal class $\Pi_i^{\mathsf{P}}$: $\mathsf{P} \subseteq
\Pi_i^{\mathsf{P}}$.
\end{theorem}

\begin{theorem}[NP lies in the first existential hierarchy level]
\lean{Complexity.NP_subset_SigmaP}
\label{Complexity.NP_subset_SigmaP}
\leanok
\uses{Complexity.NP, Complexity.SigmaP, Complexity.SigmaP_mono, Complexity.SigmaP_one}
\difficulty{1}
For every natural number $i$ with $1\le i$, every language in $\mathrm{NP}$ belongs to
the $i$th existential level $\Sigma_i^P$ of the polynomial hierarchy:
$\mathrm{NP}\subseteq\Sigma_i^P$.
\end{theorem}

\begin{theorem}[CoNP is contained in the universal hierarchy level]
\lean{Complexity.coNP_subset_PiP}
\label{Complexity.coNP_subset_PiP}
\leanok
\uses{Complexity.PiP, Complexity.PiP_mono, Complexity.PiP_one, Complexity.coNP}
\difficulty{1}
For every natural number $i$ with $1 \le i$, every language over the binary alphabet
whose complement has polynomially bounded existential certificates verifiable in
polynomial time belongs to $\Pi_i^{\mathrm P}$, the class of languages whose complements
lie in $\Sigma_i^{\mathrm P}$.
\end{theorem}

\begin{theorem}[Characterization of the polynomial hierarchy]
\lean{Complexity.mem_PH_iff}
\statementsource{Yap83}{pdf-1a201b302069-p003-b003}
\label{Complexity.mem_PH_iff}
\leanok
\uses{Complexity.PH, Complexity.SigmaP, Complexity.SigmaP_subset_SigmaP_succ}
\difficulty{4}
For every language $L$ over the binary alphabet, $L$ belongs to the polynomial hierarchy
if and only if there is a natural number $i$ such that $L\in\Sigma_i^P$.
\end{theorem}

\begin{theorem}[Existential levels lie in the polynomial hierarchy]
\lean{Complexity.SigmaP_subset_PH}
\label{Complexity.SigmaP_subset_PH}
\leanok
\uses{Complexity.PH, Complexity.SigmaP, Complexity.mem_PH_iff}
\difficulty{1}
For every natural number $i$, every language in $\Sigma_i^P$ belongs to the polynomial
hierarchy $\mathsf{PH}=\bigcup_{r\in\mathbb{N}}\Sigma_{r+1}^P$.
\end{theorem}

\begin{theorem}[Containment in the Polynomial Hierarchy]
\lean{Complexity.PiP_subset_PH}
\label{Complexity.PiP_subset_PH}
\leanok
\uses{Complexity.PH, Complexity.PiP, Complexity.PiP_subset_SigmaP_succ, Complexity.SigmaP_subset_PH}
\difficulty{1}
For every natural number $i$, the class $\Pi_i^P$ is contained in the polynomial
hierarchy $\mathrm{PH}$.
\end{theorem}

\begin{theorem}[Polynomial time lies in the polynomial hierarchy]
\lean{Complexity.P_subset_PH}
\label{Complexity.P_subset_PH}
\leanok
\uses{Complexity.P, Complexity.PH, Complexity.SigmaP_subset_PH, Complexity.SigmaP_zero}
\difficulty{1}
Every binary language decidable in polynomial time belongs to the polynomial hierarchy;
equivalently, $\mathsf{P} \subseteq \mathsf{PH}$.
\end{theorem}

\begin{theorem}[Containment of NP in the polynomial hierarchy]
\lean{Complexity.NP_subset_PH}
\label{Complexity.NP_subset_PH}
\leanok
\uses{Complexity.NP, Complexity.PH, Complexity.SigmaP_one, Complexity.SigmaP_subset_PH}
\difficulty{1}
Every language in $\mathrm{NP}$ belongs to the polynomial hierarchy $\mathrm{PH}$; that
is, $\mathrm{NP}\subseteq\mathrm{PH}$.
\end{theorem}

\begin{theorem}[CoNP containment in the polynomial hierarchy]
\lean{Complexity.coNP_subset_PH}
\label{Complexity.coNP_subset_PH}
\leanok
\uses{Complexity.PH, Complexity.PiP_one, Complexity.PiP_subset_PH, Complexity.coNP}
\difficulty{1}
Every language over the binary alphabet whose complement belongs to the class of
languages having polynomial-length existential certificates verifiable in polynomial
time belongs to the polynomial hierarchy.
\end{theorem}
